\documentclass[10pt]{article}
\usepackage[margin=1in]{geometry}
\usepackage[T1]{fontenc}
\usepackage{lmodern,microtype,amsmath,amssymb,amsthm}
\theoremstyle{definition}
\newtheorem{assumption}{Assumption}
\begin{document}
\section*{Two-Sample Quadratic Perceptron}
We study strict separation of two labeled vectors $b_+,b_-\in\mathbb{R}^d$ by a quadratic perceptron with score $F(b;\theta)$. The parameter $\theta$ belongs to $\mathbb{R}^d\times\mathbb{R}$ for the two-layer model or $\mathbb{R}^2\times\mathbb{R}^d$ for the convolutional model. The data have norms at most $R>0$ and satisfy model-specific geometric and spectral bounds; a given initialization $\theta^0$ satisfies signed wedge, depth, and size conditions. The oracle returns both signed scores and their parameter-gradient directions exactly. A fixed positive-step perceptron updates the first misclassified sample, prioritizing $b_+$. The goal is an output $\bar{\theta}$ with $F(b_+;\bar{\theta})>0$ and $F(b_-;\bar{\theta})<0$ for every realization.

Let $d\in\mathbb{N}_0$, with Euclidean inner product and norm on the feature space $\mathbb{R}^d$. Write $\mathbf{1}\in\mathbb{R}^d$ for the all-ones vector and $S$ for the cyclic shift, so that $(Sb)_i=b_{i-1\bmod d}$ when $d>0$; in dimension zero, $S$ is the unique map. The two scores are
\[
\begin{aligned}
F(b;(w,a))&=a\langle w,b\rangle
&&\text{(two-layer)},\\
F(b;((u,v),w))&=u\langle b,w\rangle+v\langle Sb,w\rangle
&&\text{(convolutional)}.
\end{aligned}
\]
Their parameter-gradient directions, in the same coordinate order, are
\[
\begin{aligned}
G(b;(w,a))&=(ab,\langle w,b\rangle),\\
G(b;((u,v),w))&=((\langle b,w\rangle,\langle Sb,w\rangle),ub+vSb).
\end{aligned}
\]
The algorithm receives only the model, $d$, $R$, and $\theta^0$. The remaining parameters introduced below may enter the oracle-call budget but are not algorithm inputs.

Our assumptions:
\begin{assumption}[Data geometry]
Both samples satisfy $\|b_+\|\leq R$ and $\|b_-\|\leq R$.
For the two-layer model, define
\[
\begin{aligned}
e_A&=\frac{b_+}{\|b_+\|},&
n&=b_--\langle b_-,e_A\rangle e_A,&
e_B&=\frac{n}{\|n\|}.
\end{aligned}
\]
We require $\langle b_+,b_-\rangle>0$ and $n\neq0$.
For the convolutional model, define the mean level and the centered directions by
\[
\begin{aligned}
\ell(b)&=\frac{\langle b,\mathbf{1}\rangle}{d},&
\widetilde b_-&=b_--\ell(b_-)\mathbf{1},\\
h_+&=\widetilde b_-+S\widetilde b_-,&
h_-&=\widetilde b_--S\widetilde b_-.
\end{aligned}
\]
We require
\[
b_+=\ell(b_+)\mathbf{1},\qquad
\ell(b_+)>0,\qquad \ell(b_-)>0,\qquad h_+\neq0.
\]
The admissible data necessarily have $d>0$; all mean levels above are taken in this case.
\end{assumption}

\begin{assumption}[Spectral bounds]
Let $g>0$, $0<c_-\leq c_+$, and $s_->0$. Define the model-dependent quantities
\[
(\lambda_1,\lambda_2,c,s)=
\begin{cases}
(\|b_+\|,\langle b_-,e_A\rangle,\|n\|,0),
&\text{two-layer},\\
(\ell(b_+)\sqrt{2d},\ell(b_-)\sqrt{2d},
\|h_+\|/\sqrt{2},\|h_-\|/\sqrt{2}),
&\text{convolutional}.
\end{cases}
\]
They satisfy
\[
\lambda_1\geq g,\qquad \lambda_2\geq g,\qquad
c_-\leq c\leq c_+.
\]
For the convolutional model, additionally $s\geq s_-$.
\end{assumption}

\begin{assumption}[Initialization geometry]
Let $\rho\geq1$ and $\omega,\delta,M>0$.
For the given starting point, define coordinates $(p,q,z,r,y)$ as follows. In the two-layer model, with $\theta^0=(w^0,a^0)$, set
\[
p=a^0,\qquad q=\langle w^0,e_A\rangle,\qquad
z=\langle w^0,e_B\rangle,\qquad r=y=0.
\]
In the convolutional model, with $\theta^0=((u^0,v^0),w^0)$, set
\[
\begin{aligned}
p&=\frac{u^0+v^0}{\sqrt{2}},&
q&=\frac{\langle w^0,\mathbf{1}\rangle}{\sqrt d},&
z&=\frac{\langle w^0,h_+\rangle}{\|h_+\|},\\
r&=\frac{u^0-v^0}{\sqrt{2}},&
y&=\frac{\langle w^0,h_-\rangle}{\|h_-\|}.
\end{aligned}
\]
Writing $\sigma=-\operatorname{sgn}(z)$, we require
\[
\begin{aligned}
|r+y|&\leq\rho|r-y|,& |z|&\geq\delta,\\
\sigma(p+q)&\geq\omega,& \sigma(p-q)&\geq0,\\
2\sigma p+2|z|&\leq M.
\end{aligned}
\]
\end{assumption}

Our result is algorithm-specific: we analyze the two-sample quadratic perceptron with a constant step size $\gamma>0$. At any query point $\theta$, one oracle call returns the tuple
\[
\bigl(F(b_+;\theta),\ G(b_+;\theta),\
-F(b_-;\theta),\ -G(b_-;\theta)\bigr).
\]
Starting at the supplied $\theta^0$, the method follows
\[
\theta^{k+1}=
\begin{cases}
\theta^k+\gamma G(b_+;\theta^k),
& F(b_+;\theta^k)\leq0,\\
\theta^k-\gamma G(b_-;\theta^k),
& F(b_+;\theta^k)>0,\ F(b_-;\theta^k)\geq0,\\
\theta^k,
& F(b_+;\theta^k)>0,\ F(b_-;\theta^k)<0.
\end{cases}
\]
The step size is fixed throughout the run and may depend only on the available inputs. After $K\in\mathbb{N}_0$ oracle calls, the output is the current iterate $\bar{\theta}=\theta^K$. The required accuracy is strict separation,
\[
F(b_+;\bar{\theta})>0,
\qquad
F(b_-;\bar{\theta})<0,
\]
for every realization of any algorithm or oracle randomness.
\end{document}
